\documentclass[12pt]{article}
\usepackage{amsmath,amssymb,amsfonts,mathtools,geometry}
\geometry{margin=1in}

\begin{document}

\title{Determinant of the Bordered Hessian (with 3$\times$3 Example)}
\author{}
\date{}
\maketitle

\paragraph{Setup.}
Let \(f:\mathbb{R}^n\to\mathbb{R}\) be maximized (or minimized) subject to \(m\) equality constraints \(g(x)=0\), where \(g:\mathbb{R}^n\to\mathbb{R}^m\).
The Lagrangian is
\[
\mathcal{L}(x,\lambda)=f(x)-\lambda^\top g(x), \qquad
g(x)=\begin{bmatrix}g_1(x)\\ \vdots \\ g_m(x)\end{bmatrix}.
\]
At a stationary point \((x^\ast,\lambda^\ast)\), define the constraint Jacobian and the Lagrangian Hessian
\[
G \;=\; \nabla g(x^\ast)\in\mathbb{R}^{m\times n},
\qquad
H \;=\; \nabla_{xx}^2 \mathcal{L}(x^\ast,\lambda^\ast)\in\mathbb{R}^{n\times n}.
\]
The \emph{bordered Hessian} is the \((m+n)\times(m+n)\) block matrix
\[
B \;=\;
\begin{bmatrix}
\mathbf{0}_{m\times m} & G\\
G^\top & H
\end{bmatrix}.
\]

\paragraph{Determinant via block/Schur complement.}
Assume \(H\) is invertible. Using the block determinant identity,
\[
\det\!\begin{bmatrix}A&B\\C&D\end{bmatrix}
=\det(D)\,\det(A-BD^{-1}C),
\]
with \(A=\mathbf{0},\,B=G,\,C=G^\top,\,D=H\), we get
\[
\boxed{\;
\det(B)\;=\;\det(H)\,\det\!\bigl(\,A-BD^{-1}C\,\bigr)
\;=\;\det(H)\,\det\!\bigl(-\,G\,H^{-1}G^\top\bigr)
\;=\;(-1)^m\,\det(H)\,\det\!\bigl(GH^{-1}G^\top\bigr).
\;}
\]
For \(m=1\) (one constraint), writing \(G=[\,g_1\;\cdots\;g_n\,]\),
\[
\boxed{\;\det(B)\;=\;-\;\det(H)\;\;G\,H^{-1}G^\top.\;}
\]

\paragraph{Exact \(3\times 3\) bordered Hessian (one constraint, two variables).}
Take \(n=2,\; m=1\). Write
\[
H=\begin{bmatrix} a & b\\ b & c\end{bmatrix},\qquad
G=\begin{bmatrix} g_1 & g_2\end{bmatrix}.
\]
Then the bordered Hessian is
\[
B \;=\;
\begin{bmatrix}
0 & g_1 & g_2 \\
g_1 & a & b \\
g_2 & b & c
\end{bmatrix}.
\]
The \(2\times 2\) Hessian determinant is \(\Delta=\det(H)=ac-b^2\), and
\[
H^{-1}=\frac{1}{\Delta}\begin{bmatrix} c & -b\\ -b & a\end{bmatrix}.
\]
Hence
\[
G\,H^{-1}G^\top
= \frac{1}{\Delta}\bigl(c\,g_1^2 - 2b\,g_1 g_2 + a\,g_2^2\bigr),
\qquad
\det(B) \;=\; -\,\det(H)\cdot G H^{-1}G^\top
\;=\; -\bigl(c\,g_1^2 - 2b\,g_1 g_2 + a\,g_2^2\bigr).
\]

\paragraph{Direct hand expansion check (same \(3\times 3\)).}
Expanding \(\det(B)\) along the first row:
\[
\det(B)
= 0\cdot(\cdots)
- g_1\cdot \det\!\begin{bmatrix} g_1 & b\\ g_2 & c\end{bmatrix}
+ g_2\cdot \det\!\begin{bmatrix} g_1 & a\\ g_2 & b\end{bmatrix}.
\]
Compute the \(2\times 2\) minors:
\[
\det\!\begin{bmatrix} g_1 & b\\ g_2 & c\end{bmatrix}=g_1 c - b g_2,
\qquad
\det\!\begin{bmatrix} g_1 & a\\ g_2 & b\end{bmatrix}=g_1 b - a g_2.
\]
Therefore
\[
\det(B)
= -\,g_1(g_1 c - b g_2)\;+\;g_2(g_1 b - a g_2)
= -\,c\,g_1^2 + 2b\,g_1 g_2 - a\,g_2^2
= -\bigl(c\,g_1^2 - 2b\,g_1 g_2 + a\,g_2^2\bigr),
\]
which exactly matches the Schur complement result above.

\paragraph{Bordered principal minors and second-order tests (quick reference).}
For \(m\) equality constraints, define the bordered principal minors of order \(m+k\) (border first, then successively add one decision variable):
\[
\Delta_{m+k} \;=\; \det\bigl(\text{leading }(m+k)\times(m+k)\text{ principal submatrix of }B\bigr),
\quad k=1,\ldots,n.
\]
\begin{itemize}
\item \textbf{Maximization:} At a regular stationary point, a sufficient condition for a local maximum is
\[
(-1)^m\,\Delta_{m+1}>0,\quad (-1)^m\,\Delta_{m+2}>0,\;\ldots,\;(-1)^m\,\Delta_{m+n}>0.
\]
\item \textbf{Minimization:} A sufficient condition for a local minimum is
\[
(-1)^m\,\Delta_{m+1}<0,\quad (-1)^m\,\Delta_{m+2}<0,\;\ldots,\;(-1)^m\,\Delta_{m+n}<0.
\]
\end{itemize}
These criteria use the signs of bordered minors rather than \(\det(B)\) alone. When \(H\) is invertible, the Schur formula above is a fast way to compute \(\det(B)\) (and the top-order minor) exactly.

\paragraph{Remarks.}
\begin{itemize}
\item The ordering of rows/columns matters: the standard convention borders the constraints first, then appends the Hessian block \(H\).
\item If \(H\) is singular, the Schur complement formula is not applicable as written. In that case, compute \(\det(B)\) directly (e.g., expansion or row/column operations that preserve determinants).
\item For multiple constraints, \(\det(GH^{-1}G^\top)\) is the squared volume (up to scaling) induced by the metric \(H^{-1}\) on the constraint gradients; the \((-1)^m\) factor tracks the $m$-fold sign from the zero block on the border.
\end{itemize}

\end{document}
